\scanpage{52}
\uline{Cellules}

$X$ espace topologique, $(F_i)_{i\in I}$ une partition de $X$ (les $F_i$ sont appelés ``feuilles''). On dit que la partition est \uline{admissible} si
\[
\mathrm I\left\{\begin{array}{ll}
\inserted{(0)}&\inserted{(F_i)\text{ est une partition de }X}\\
(i)&(F_i)\text{ est loc. finie }\deleted{\unsure{et recouvrent}}\\
(ii)&\forall i,\;\overline F_i\text{ est une réunion de feuilles, et}\\
(iii)&F_i\text{ ouvert dans }\overline F_i\quad(\text{i.e. }F_i\text{ loc. fermé dans }X)
\end{array}\right.
\]

Les $\overline F_i$ sont appelés ``feuilles fermées''. Quand on a $\overline F_i=\overline F_j$, on a $F_i=F_j$, car si $\overline F_j\subset\overline F_i$ i.e. $F_j\subset\overline F_i$, comme $\overline F_i$ est réunion de feuilles, on a $F_j=F_i$, ou $F_j\subset\overline F_i-F_i$, mais \deleted{\unsure{alors}} \inserted{dans ce cas} $\overline F_j\subset\overline F_i-F_i$ donc $\overline F_j\neq\overline F_i$ puisque $F_i\neq\varnothing$. Donc la donnée de la famille $(F_i)$ équivaut à celle de la famille des $\Phi_i=F_i$. Celle-ci \unsure{satisfait} :
\[
\mathrm{II}\left\{\begin{array}{ll}
(0)&\text{les }\Phi_i\;\deleted{\unsure{non vides}}\;\text{recouvrent }X\\
(i)&(\Phi_i)\text{ est loc. finie}\\
(ii)&\text{les }\Phi_i\text{ fermés}\\
(iii)&\forall i,j,\;\Phi_i\cap\Phi_j\text{ est une réunion de }\Phi_k\\
(iv)&\forall i,j\text{ avec }\Phi_j\subsetneq\Phi_i,\;
\deleted{\Phi_j\neq\Phi_i}\;\Phi_j\text{ est \uline{rare} dans }\Phi_{\unsure{i}}
\end{array}\right.
\]

Ces conditions, néc. pour que $(\Phi_i)$ soit la famille des feuilles fermées d'une partition admissible, est aussi suffisante. En effet, $\forall i\in I$, soit \deleted{$\unsure{F}$}
\[
F_i=\Phi_i-\bigcup_{\Phi_j\subsetneq\Phi_i}\Phi_j
\]
les $F_i$ forment une \uline{partition} de $X$ ($\forall x\in X$, $\exists$ unique $i$ tel que
\margin{$x\in F_i$, savoir $i$ tel que $x\in\Phi_i$ et $\Phi_i$ soit le plus petit de ceux-ci (existe grâce à (i) (iii))}
$F_i$ est loc. fermé grâce à (i) et (ii)), $(F_i)$ est loc. fini grâce à (i), \deleted{$\unsure{F_\alpha}$} $\forall i\quad\overline F_i=\Phi_i$ (grâce à (iv)),
prouvons que $\Phi_i$ est réunion de feuilles $F_j$, i.e.
\deleted{Considérons le \unsure{volet}}
\deleted{\unsure{Soit} $\unsure{\forall}$}
que $\forall j,\quad F_j\cap\Phi_i\neq\varnothing\Rightarrow F_j\subset\Phi_i$.
Mais $F_j\cap\Phi_i\neq\varnothing\Longrightarrow\deleted{\unsure{\Phi_j}}\;\Phi_j\subset\Phi_i$ (prendre $x\in F_j\cap\Phi_i$, et utiliser la description marginale de $j$ tel que $x\in F_j$), d'où $F_j\subset\Phi_i$ OK,

\margin{53}
